\begin{claim}{D.46}{Proved}{Three positive involution sections and order-four lifting}
Let $T$ be the full once-punctured-torus loop-outer group
of $F=\widehat F(x,y)$.
For every involution $b\in T$ whose $2$-adic orientation
is positive, its centre-free extension
\[
             1\longrightarrow F\longrightarrow E_b
                  \longrightarrow\langle b\rangle
                  \longrightarrow1
\]
has exactly three $F$-conjugacy classes of order-two
sections.  For every such section $S$,
\[
                         H^1(C_F(S),\F_2)=0.             \tag{D.46a}
\]
Every cyclic order-four subgroup of $T$ has a homomorphic
lift to $\operatorname{Aut}(F)$, with its prescribed
outer action.  Thus every element of order four in
$T$ has an actual order-four representative.
\end{claim}
\begin{proof}
\emph{The homology of a positive involution.}
D.27 gives an actual order-two representative $\beta$
of $b$, and identifies $E_b$ with
$F\rtimes_\beta\langle s\rangle$.
The map of D.33 on the finite $2$-group $\langle b\rangle$
is faithful.  The matrix of $\beta$ on
$H_1(F,\Z_2)$ is therefore nonidentity.
Its square is the identity and its determinant is $1$.
Over $\Q_2$ it is diagonalizable; the determinant
condition makes its two eigenvalues equal.
It follows that
\[
                    \beta^{\rm ab}_2=-I.               \tag{D.46b}
\]
This uses the punctured-torus statement D.33 to exclude
the identity matrix; no analogous exclusion for an
arbitrary unmarked rank-two action is assumed.

\emph{The full cohomology algebra and its Bockstein.}
Let $u:E_b\to\F_2$ be the quotient character, with
$u(s)=1$.  Choose a basis $\chi_1,\chi_2$ of
$H^1(F,\F_2)$ and extend these characters to $E_b$
by declaring $\chi_i(s)=0$.
These extensions exist because (D.46b) is the identity
modulo $2$.
The Hochschild--Serre spectral sequence has precisely
the two rows
\[
 E_2^{a,0}=\F_2 u^a,\qquad
 E_2^{a,1}=H^1(F,\F_2)\,u^a.                            \tag{D.46c}
\]
The transgression on $E_2^{0,1}$ is zero since the two
characters have just been extended to $E_b$.
Its other possible differentials are the products of
this transgression with powers of $u$, and hence are
zero as well.  There are no further nonzero rows.
Consequently, for every $n\geq1$,
\[
 H^n(E_b,\F_2)
   =\langle u^n,u^{n-1}\chi_1,u^{n-1}\chi_2\rangle_{\F_2},
 \qquad \dim H^n(E_b,\F_2)=3.                           \tag{D.46d}
\]
Independence follows from the two associated graded
pieces in (D.46c).
Multiplication by $u$ is an isomorphism in every
degree at least one: it is an isomorphism on those
associated graded pieces.

For each $i$, lift $\chi_i|_F$ to a continuous
homomorphism $\varphi_i:F\to\Z/4$.
Equation (D.46b) gives
$\varphi_i(\beta(h))=-\varphi_i(h)$.
The actual continuous one-cochain on $E_b$
\[
             \widetilde\chi_i(hs^\epsilon)=\varphi_i(h)
                      \quad(\epsilon=0,1)
\]
therefore satisfies
\[
 \begin{split}
 \delta\widetilde\chi_i(hs^\epsilon,ks^\eta)
  &=\varphi_i(h)+\varphi_i(k)
             -\varphi_i(h\beta^\epsilon(k))\\
  &=2\epsilon\,\varphi_i(k)
       \quad\text{in }\Z/4.
 \end{split}                                           \tag{D.46e}
\]
For the coefficient sequence
$0\to\F_2\xrightarrow{\;2\;}\Z/4\to\F_2\to0$,
this says that the Bockstein of $\chi_i$ is
$u\chi_i$.
It is also $\chi_i^2$.
For completeness, the lift of any $\F_2$ character
which takes its values in $\{0,1\}\subset\Z/4$
has coboundary, divided by $2$, equal to the cup
square of that character; comparison of two lifts
changes this cocycle only by a coboundary.
Thus
\[
                            \chi_i^2=u\chi_i.           \tag{D.46f}
\]

Invert the central class $u$ and take degree zero:
\[
             R=\bigl(H^*(E_b,\F_2)[u^{-1}]\bigr)^0.
\]
It has dimension three, with basis
$1,z_1=\chi_1/u,z_2=\chi_2/u$ by (D.46d).
Equation (D.46f) makes both $z_i$ idempotent.
In particular $R$ is a quotient of the algebra
\[
 \F_2[Z_1,Z_2]/(Z_1^2-Z_1,Z_2^2-Z_2)
           \cong\prod_{(a_1,a_2)\in\F_2^2}\F_2,
 \quad P\longmapsto(P(a_1,a_2))_{(a_1,a_2)}.             \tag{D.46g}
\]
The displayed evaluation map is an isomorphism:
the four products of $Z_i$ and $1-Z_i$ are its four
coordinate idempotents.
Every ideal of a finite product of fields is a product
of ideals in those fields.  Every quotient is therefore
reduced, so $R$ has no nonzero nilpotent element.

\emph{Brown restriction and the three full sections.}
Every finite subgroup of $E_b$ injects into
$E_b/F=C_2$, since $F$ is torsion-free.
Thus the nontrivial elementary abelian $2$-subgroups
are exactly its order-two sections, with no inclusions
between distinct ones.
For such a section $S$ put $C_S=C_F(S)$.
Its normalizer is the exact direct product
\[
                         N_{E_b}(S)=C_S\times S.        \tag{D.46h}
\]
Indeed every normalizer element centralizes $S$, since
$\operatorname{Aut}(C_2)=1$; multiplying it by its
unique image in $S$ leaves an element of $C_S$.

The full-profinite Brown restriction theorem used in
D.28 and D.33 applies because $F$ is open in $E_b$
and has $2$-cohomological dimension one
\cite[Theorem~2 and \S1.1]{A-KochloukovaZalesskii2007}.
In sufficiently high degrees it identifies cohomology
with the product of the cohomologies of (D.46h);
the identification is the actual restriction map and
preserves cup products.
Each factor contributes a copy of $\F_2$ in every
positive degree, so (D.46d) in particular makes the
number of factors finite.
Since $C_S$ is a closed subgroup of $F$, it has
$\operatorname{cd}_2(C_S)\leq1$.
After the same localization, restriction identifies
$R$ as a ring with
\[
       \prod_{[S]}
         \bigl(\F_2\oplus H^1(C_S,\F_2)\,u^{-1}\bigr).
                                                        \tag{D.46i}
\]
The second summand in each factor is an ideal with
square zero, since $H^2(C_S,\F_2)=0$.
Here $u|_{C_S}=0$ and $u|_S$ is the nonzero
character, so the localization in (D.46i) is exactly
the localization already used in (D.46g).
Reducedness of $R$ forces every second summand to
vanish.  Its dimension three then proves that there
are exactly three section classes, and proves
(D.46a).

These are also precisely the $F$-conjugacy classes.
To see this directly, let $e\in E_b$ and choose
$s'\in S$ with the same image as $e$ in $E_b/F$.
Then $eS e^{-1}=(es'^{-1})S(es'^{-1})^{-1}$,
and $es'^{-1}\in F$.

\emph{An invariant section in the order-four extension.}
Now let $a\in T$ have order four, and form its
centre-free extension
\[
 E=\{(\phi,j)\in\operatorname{Aut}(F)\times C_4:
                                      [\phi]=a^j\}.
                                                        \tag{D.46j}
\]
Its kernel is identified with $F$ by
$h\mapsto(\operatorname{Inn}(h),0)$.
Let $E_2$ be the inverse image of $\langle a^2\rangle$.
The orientation of $a^2$ at $2$ is positive:
the only roots of unity of order dividing four
in $\Z_2^\times$ are $1,-1$, whose squares are $1$.
Equivalently, D.33 embeds $\langle a\rangle$ in
$\operatorname{GL}_2(\Z_2)$, where its square is $-I$.
The first part of the proof thus applies to $E_2$.

Conjugation by $E$ permutes the three $F$-conjugacy
classes of sections of $E_2$.
Every element of $E_2$ acts trivially on this class
set, by the preceding paragraph, so the permutation
action factors through $E/E_2=C_2$.
An action of $C_2$ on a set of three elements has a
fixed element.
Choose a section $S$ in that class.
If $e\in E$ projects to the specified generator of
$C_4$, there is $f\in F$ with
$eSe^{-1}=fSf^{-1}$.
Hence $f^{-1}e\in N_E(S)$ still projects to that
generator.
Put $N=N_E(S)$ and $C=C_F(S)$.
We have proved the exact sequences
\[
 \begin{gathered}
 1\longrightarrow C\longrightarrow N
              \longrightarrow C_4\longrightarrow1,
       \qquad S\subseteq Z(N),\\
 1\longrightarrow C\longrightarrow J:=N/S
              \longrightarrow C_2\longrightarrow1.
 \end{gathered}                                         \tag{D.46k}
\]
The embedding of $C$ in the second line is injective
because $C\cap S=1$.

\emph{Projectivity supplies the order-four section.}
Equation (D.46a) implies
\[
                              C^{(2)}=1:                \tag{D.46l}
\]
every nontrivial pro-$2$ group has a nontrivial finite
$2$-group quotient, which in turn has a quotient
$C_2$.
In particular every homomorphism from $J$ to a finite
$2$-group kills $C$ and factors through the last
$C_2$ in (D.46k).  That quotient is nontrivial, so
\[
                              J^{(2)}=C_2.              \tag{D.46m}
\]

We show that $J$ has a nonidentity torsion element.
If it were torsion-free, its open subgroup $C$ has
$\operatorname{cd}_p(C)\leq1$ at every prime $p$,
as a closed subgroup of the free profinite group $F$.
Serre's torsion-free finite-index dimension theorem,
in exactly the form used in D.27, would give
$\operatorname{cd}_p(J)\leq1$ at every prime.
Thus $J$ would be projective as a full profinite
group.
Its maximal pro-$2$ quotient would then be a
projective, hence free, pro-$2$ group: pull a finite
$2$-group embedding problem back to $J$, solve it
by projectivity, and factor the solution through
$J^{(2)}$.
This contradicts (D.46m), since $C_2$ is not a free
pro-$2$ group.

Let $j\in J$ be a nonidentity torsion element.
Its cyclic subgroup injects into $J/C=C_2$, because
$C$ is torsion-free.
Consequently $j$ has order two and lies outside $C$.
Choose $n\in N$ over $j$.
Then $n^2\in S$, whereas its image in $C_4$ is an
odd element and has order four.
It follows that $n^2$ is the nonidentity element
of $S$, and $n$ has order exactly four.
After replacing $n$ by $n^{-1}$ if necessary, its
projection is the specified generator of $C_4$.
Thus $C_4\to E$, sending that generator to $n$,
is a section of (D.46j).
Projection to $\operatorname{Aut}(F)$ gives the
claimed homomorphic lift with its original labels.

The program \texttt{code/part\_c/torus\_order4\_lift\_d46.py}
checks the finite cochain formula (D.46e), its
comparison with the cup square, and the explicit
Boolean-algebra evaluation in (D.46g).
The dimension and normalizer arguments above prove
the statement for the full profinite extensions.
No claim of conjugacy to a geometric order-four
action is made.  D.43 applied to $a^2$ does also
give $\lambda(a)^2=1$ globally, so the remaining
orientation issue for $a$ concerns mixed signs,
rather than fourth roots of unity.
\end{proof}
